\section{A positive shrinking profile}\label{prof:section}
\label{sec:profile}

We construct the profile used in the concentration argument.  Its construction
has two parts: a finite shooting estimate, proved by rational enclosures in
Appendix~\ref{cert:section}, and an analytic construction on a half-line.
The finite calculation belongs to the established framework of
validated Taylor integration for ordinary differential
equations~\cite{NedialkovJacksonCorliss1999}. Here the integer interval
recurrences and complex-disk remainder bounds are given explicitly.
The half-line solution is constructed by a contraction, whereas the
final matching argument uses only continuity of the shooting map.

\begin{theorem}[Shrinking profile]\label{prof:profile}
There exist a constant $c>0$ and a function $G\in C^\infty([0,\infty))$
such that
\begin{equation}\label{prof:profile-equation}
 Q(y\partial_y)(G-1)=\frac{y^2}{2}\left(\frac1G-\frac1c\right),
 \qquad Q(D)=(D-1)(D+9)(D+10).
\end{equation}
The function $G$ is bounded above and bounded away from zero, satisfies
$G(0)=1$ and $G'(0)<0$, and is analytic near zero.  There are real
coefficients $g_j$, $j\ge1$, such that, for every $J,\ell\ge0$,
\begin{equation}\label{prof:expansion}
 \frac{\dd^\ell}{\dd y^\ell}
 \left(G(y)-c-\sum_{j=1}^{J}g_jy^{-2j}\right)
 =O\bigl(y^{-2(J+1)-\ell}\bigr)\qquad(y\longrightarrow\infty).
\end{equation}
In particular, $u\mapsto G(u^{-1/2})$ extends smoothly to $u=0$ from
the right, with value $c$.

Set $h(y)=(G(y)-1)/y$, with its smooth value at zero.  For $t<0$, the
function
\begin{equation}\label{prof:flat-shrinker}
 F(t,x)=\rho^{-1}h(x/\rho),\qquad \rho=\sqrt{-t},
\end{equation}
solves
\[
 F_t+(1+xF)^2\widehat{\A}F=0,
 \qquad \widehat{\A}=(x\partial_x^2+12\partial_x)^2.
\]
\end{theorem}

\subsection{Reduction and the analytic family at the origin}

The calculation in Section~\ref{geo:core-section} shows that
\eqref{prof:profile-equation} makes $F$ in
\eqref{prof:flat-shrinker} an exact solution of the flat equation.
It remains to construct a positive $G$ with the stated behavior at
zero and infinity.

Put
\begin{equation}\label{prof:change-variables}
 c=(1+d)^{-1},\qquad
 s=\left(\frac{27}{16}\right)^{1/3}\left(\frac yc\right)^{2/3},
 \qquad G(y)=c(1+p(s)).
\end{equation}
Since $y\partial_y=\frac23s\partial_s$ and $G-1=c(p-d)$,
\eqref{prof:profile-equation} becomes
\begin{equation}\label{prof:ode}
 \mathcal E_d(p):=
 p'''+\frac{30}{s}p''+\frac B{s^2}p'
       +\frac C{s^3}(p-d)+\frac p{1+p}=0,
 \qquad B=\frac{751}{4},\quad C=-\frac{1215}{4}.
\end{equation}
Primes in this section henceforth denote $s$-derivatives, except where
another variable is explicitly displayed.  For example, the identity
$Q(D)=D^3+18D^2+71D-90$ verifies all the coefficients in
\eqref{prof:ode}.

\Needspace{4\baselineskip}
We will construct origin solutions
\(p(s)=d+ks^{3/2}+O(s^3)\), with two real parameters: \(d\) fixes
\(c=(1+d)^{-1}\), and \(k\) fixes \(G'(0)\).  These parameters must
be chosen to match \(p\) to a tail tending to zero.  The limiting
linearized equation \(p'''+p=0\) has roots \(-1\) and
\((1\pm i\sqrt3)/2\): one decaying direction and two real growing
directions.

The tail argument uses a constant-coefficient comparison equation
with the same splitting.  After subtracting an explicit decaying
approximation, we express the residual jet at \(S_0=56\) in its
eigenbasis.  For each allowed \(d\), a contraction constructs small
bounded tails with one real coordinate free and the other two
determined.  Choosing the free coordinate from the origin solution leaves two
real matching conditions to enforce by varying \((k,d)\).

The affine square below makes the shot's two remaining coordinates
close to its parameter \(z_*\).  Together with the tail bound, the
certified estimate yields a continuous self-map of the square;
Brouwer's theorem supplies a matching point.  All terminating
decimals in this section denote exact rational numbers.  Define
\begin{equation}\label{prof:parameters}
 \begin{gathered}
 R=5\cdot10^{-10},\qquad
 \mathcal Z=\{z_*\in\R^2:\norm{z_*}_\infty\le R\},\qquad S_0=56,\\
 \begin{pmatrix}k\\d\end{pmatrix}
 =\begin{pmatrix}-.094617116666041\\ .400682721856639\end{pmatrix}
 +10^{-8}
 \begin{pmatrix}25190&-63731\\721855&359157\end{pmatrix}z_*.
 \end{gathered}
\end{equation}
Write $k_0,d_0$ for the central values.  The square implies
\begin{equation}\label{prof:parameter-displacement}
 |k-k_0|\le4.44605\cdot10^{-13},\qquad
 |d-d_0|\le5.40506\cdot10^{-12}<.011R.
\end{equation}

\begin{lemma}[Origin family]\label{prof:origin}
For $|k-k_0|\le .1$ and $|d-d_0|\le .1$ in $\C$, there is a unique
solution germ of \eqref{prof:ode}, analytic in $r=s^{3/2}$, of the form
\begin{equation}\label{prof:origin-series}
 p=d+kr+\sum_{i\ge2}a_ir^i,\qquad r=s^{3/2}.
\end{equation}
The series is jointly holomorphic in $r,k,d$ for $|r|<3$ and satisfies
\begin{equation}\label{prof:origin-majorant}
 3|k|+\sum_{i\ge2}|a_i|3^i<.8.
\end{equation}
Each of $p(1),p'(1),p''(1)$ has modulus at most $2$ on this parameter
polydisc.  For the real parameters in \eqref{prof:parameters}, $1+p>0$
on $0\le s\le1$.
\end{lemma}

\begin{proof}
In the variable $r=s^{3/2}$, \eqref{prof:ode} is equivalent to
\[
 Q(r\partial_r)(p-d)=\frac8{27}r^2\bigl((1+p)^{-1}-1\bigr).
\]
Consequently the coefficients of an analytic solution must satisfy
\begin{equation}\label{prof:origin-recurrence}
 a_i=
 \frac{[r^{i-2}]\bigl((1+p)^{-1}-1\bigr)}
 {(27/8)(i-1)(i+9)(i+10)},\qquad i\ge2.
\end{equation}
The constant and linear equations leave $d,k$ free, and this recurrence
determines every subsequent coefficient uniquely.

For convergence, use the Banach algebra of power series with norm
$\sum_i|b_i|3^i$.  On the closed ball of series $a$ with no constant or
linear term and norm at most $.2$, one has
\[
 \norm{kr+a}<.585+.2<.8,\qquad |1+d|>1.30.
\]
Thus inversion of $1+d+kr+a$ is given by a convergent geometric series
in that Banach algebra and has norm at most $2$.  The map specified by
\eqref{prof:origin-recurrence} sends this ball into the ball of radius
$9(2+1)/440<.062$, since every denominator in that formula exceeds $440$.
Its Lipschitz constant is at most $9\cdot 2^2/440<1$.  Iteration therefore
gives a unique fixed point, jointly holomorphic in the parameters.  The
strict inequalities also hold on a neighborhood of the closed parameter
polydisc.  Moreover
\[
 3|k|+\norm{a}
 \le 3(.094617116666041+.1)+\frac{27}{440}<.8,
\]
which proves \eqref{prof:origin-majorant}.

For $0\le r\le1$, the nonconstant sum is at most $.8r/3$.  In particular
the real parameters in \eqref{prof:parameters} give $1+p>1.4-.8/3>0$.
At $s=1$, a term $a_is^{3i/2}$ acquires factors $3i/2$ and
$(3i/2)(3i/2-1)$ under one and two derivatives.  The elementary bounds
\[
 \sup_{i\ge1}\frac{3i/2}{3^i}=\frac12,
 \qquad
 \sup_{i\ge1}\frac{(3i/2)(3i/2-1)}{3^i}=\frac23
\]
together with \eqref{prof:origin-majorant} and $|d|<.501$ give the stated
jet bounds.  Termwise differentiation is valid since $1<3$.
\end{proof}

\subsection{The matching coordinates}

Define the four-term approximation
\begin{equation}\label{prof:asymptotic-polynomial}
 \begin{gathered}
 p_{\rm a}(s;d)=\sum_{j=1}^4d_js^{-3j},\qquad
 a=d_1=Cd,\qquad b=d_2=a^2+567a,\\
 d_3=2ab-a^3+\frac{2025}{4}b,\qquad
 d_4=b^2+2ad_3-3a^2b+a^4+\frac{567}{2}d_3.
 \end{gathered}
\end{equation}
We choose the constant coefficients in $D_0$ at the midpoints of the
corresponding linearized tail ranges on $[S_0,\infty)$, giving the
uniform difference bounds used in Lemma~\ref{prof:tail}.  Let
\begin{equation}\label{prof:coordinates}
 \begin{gathered}
 D_0(\lambda)=\lambda^3+\frac{15}{S_0}\lambda^2
       +\frac{B}{2S_0^2}\lambda+1+\frac{C}{2S_0^3},\\
 v_i=(1,\lambda_i,\lambda_i^2)^t,\qquad
 \begin{pmatrix}p-p_{\rm a}\\p'-p_{\rm a}'\\p''-p_{\rm a}''\end{pmatrix}_{s=S_0}
       =qv_0+\xi v_1+\bar\xi v_2.
 \end{gathered}
\end{equation}
The roots are ordered so that $\lambda_0$ is real,
$\Im\lambda_1>0$, and $\lambda_2=\bar\lambda_1$.  The exact real
projection matrix $\Gamma$ takes a real jet to $(q,\Re\xi,\Im\xi)^t$.
Its rational approximation used in the finite certificate is
\begin{equation}\label{prof:projection-approximation}
 \Lambda=10^{-5}
 \begin{pmatrix}
 30756&-27384&33447\\
 34622&13692&-16724\\
 -2925&-34040&-28850
 \end{pmatrix}.
\end{equation}

\begin{lemma}[Root and convolution bounds]\label{prof:root-bounds}
The roots in \eqref{prof:coordinates} are distinct and satisfy
\begin{equation}\label{prof:root-boxes}
 \begin{split}
 -1.08659&<\lambda_0<-1.08655,\\
 .40933&<\Re\lambda_1<.40938,\qquad
 .8671<\Im\lambda_1<.86721.
 \end{split}
\end{equation}
Every entry of $\Gamma-\Lambda$ has modulus less than $.001$.
The absolute row sums of $\Gamma$ are less than $.92,.68,.68$,
respectively.  With
\begin{equation}\label{prof:convolution-constants}
 M_i=\sum_{j=0}^2
 \frac{|\lambda_j|^i}{|\Re\lambda_j|\,|D_0'(\lambda_j)|},\qquad 0\le i\le2,
\end{equation}
one has
\[
 M_0<1.96,\qquad M_1<1.92,\qquad M_2<1.88,
 \qquad \frac1{\Re\lambda_1|D_0'(\lambda_1)|}<.82.
\]
\end{lemma}

\begin{proof}
Write $\alpha=15/56$, $\beta=B/(2\cdot56^2)$ and
$H_0=1+C/(2\cdot56^3)$.  The exact calculations
\[
 \alpha^2-3\beta=-\frac{453}{25088}<0,\qquad
 D_0(-1.08659)=-\frac{17566501197647}{343000000000000000}<0,
\]
\[
 D_0(-1.08655)=\frac{187622583137}{2744000000000000}>0
\]
show that $D_0$ is strictly increasing on the real line and has exactly
one real root in the indicated interval.  Write $\lambda_0=-r_0$ and
$\lambda_1=u+iv$ with $v>0$.  Vieta's identities give
\begin{equation}\label{prof:root-arithmetic}
 u=\frac{r_0-\alpha}{2},\qquad
 v^2=\frac{H_0}{r_0}-u^2,\qquad |\lambda_1|^2=\frac{H_0}{r_0}.
\end{equation}
In particular, if $u_-=114617/280000$ and $u_+=573113/1400000$, then
$u_-<u<u_+$, with $.40933<u_-$ and $u_+<.40938$.  The rational inequalities
\[
 .8671^2<\frac{H_0}{1.08659}-u_+^2
 \le v^2\le\frac{H_0}{1.08655}-u_-^2<.86721^2,
 \qquad \frac{H_0}{1.08655}<.959^2
\]
give the other root bounds and $|\lambda_1|<.959$.

For each root, the corresponding complex projection row is
\begin{equation}\label{prof:projection-row}
 \frac{(\lambda_j\lambda_k,-\lambda_j-\lambda_k,1)}
 {D_0'(\lambda_i)},\qquad \{i,j,k\}=\{0,1,2\}.
\end{equation}
Taking the real row for $i=0$ and the real and imaginary parts for $i=1$
gives $\Gamma$.  Put $E=(r_0+u)^2+v^2$.  The three real rows are
\[
 \begin{gathered}
 \left(\frac{u^2+v^2}{E},\frac{-2u}{E},\frac1E\right),\\
 \left(\frac{r_0(r_0+2u)}{2E},\frac uE,\frac{-1}{2E}\right),\\
 \left(\frac{r_0[u(r_0+u)-v^2]}{2vE},
       \frac{-(v^2+r_0^2-u^2)}{2vE},\frac{-(r_0+u)}{2vE}\right).
 \end{gathered}
\]
Substituting $r_0\in[1.08655,1.08659]$ and the intervals just obtained
for $u,v,v^2$ gives the following rational enclosures.  For each entry,
products are enclosed by the minimum and maximum of the four endpoint
products, and positive reciprocals by reversing the reciprocal endpoints.
\begin{center}
\begin{tabular}{c c c c}
\toprule
 & First entry & Second entry & Third entry\\
\midrule
$q$ & $(.3074,.3078)$ & $(-.2740,-.2737)$ & $(.3344,.3346)$\\
$\Re\xi$ & $(.3461,.3464)$ & $(.1368,.1370)$ & $(-.1673,-.1672)$\\
$\Im\xi$ & $(-.0294,-.0291)$ & $(-.3406,-.3402)$ & $(-.2887,-.2883)$\\
\bottomrule
\end{tabular}
\end{center}
The entrywise error from $\Lambda$ is therefore less than $.001$, and
the row sums are less than $.9164,.6507,.6587$, respectively.

Finally,
\[
 E>(1.08655+u_-)^2+.8671^2
   =\frac{46876434629}{15680000000}>2.9895>1.729^2.
\]
Since $D_0'(-r_0)=E$ and $|D_0'(u+iv)|=2v\sqrt E$, for $i=0,1,2$ one has
\[
 M_i<\frac{1.08659^i}{1.08655\cdot2.9895}
      +\frac{.959^i}{.40933\cdot.8671\cdot1.729}.
\]
These three rational expressions are less than $1.96,1.92,1.88$,
respectively, by multiplication by their positive denominators.  Similarly,
\[
 \frac1{\Re\lambda_1|D_0'(\lambda_1)|}
 <\frac1{2\cdot.40933\cdot.8671\cdot1.729}
 =\frac{500000000000}{613674044347}<.82.
\]
\end{proof}

The finite part of the construction is the following consequence of
Proposition~\ref{cert:shoot}.  The origin solution of
Lemma~\ref{prof:origin}, for every $z_*\in\mathcal Z$, continues to
$s=S_0$ and satisfies
\begin{equation}\label{prof:forward-bounds}
 1+p(s)>0\quad(0\le s\le S_0),\qquad
 |q_{\rm shot}(z_*)|\le .25R,
 \qquad
 \norm{(\Re\xi_{\rm shot},\Im\xi_{\rm shot})-z_*}_\infty\le .13R.
\end{equation}
The proof in the Appendix includes the analytic interpretation of its
rational Taylor enclosures, the propagation of all rounding and truncation
errors, and the uniform nonlinear continuation estimate.  In particular,
\eqref{prof:forward-bounds} is a statement about the exact solutions.
Their endpoint jets depend continuously on $z_*$: this follows first from
Lemma~\ref{prof:origin} at $s=1$, and then from ordinary continuous
dependence for \eqref{prof:ode} where $s>0$ and $1+p>0$.
The estimates from this section used in Appendix~\ref{sec:certificate}
are the origin family and the explicit root and tail-polynomial estimates;
all are independent of the matching step.

\subsection{Construction on the half-line}

\begin{lemma}[Tail residual]\label{prof:tail-residual}
Uniformly for the parameters in \eqref{prof:parameters},
\[
 \sup_{s\ge56}|\mathcal E_d(p_{\rm a})(s)|<3\cdot10^{-13}.
\]
Each derivative of this residual tends to zero at infinity.
If $P_{\rm a}(d)=(p_{\rm a},p_{\rm a}',p_{\rm a}'')(56;d)$, then
\begin{equation}\label{prof:tail-parameter-bound}
 \norm{P_{\rm a}(d)-P_{\rm a}(d_0)}_\infty\le .01|d-d_0|.
\end{equation}
\end{lemma}

\begin{proof}
Put $t=s^{-3}$.  For a monomial $d_js^{-3j}$, the differential part
$p'''+30s^{-1}p''+Bs^{-2}p'+Cs^{-3}p$ contributes
$L(-3j)d_jt^{j+1}$, where
\[
 L(\nu)=\nu(\nu-1)(\nu-2)+30\nu(\nu-1)+B\nu+C.
\]
Its first four relevant values are
\[
 L(-3)=-567,\quad L(-6)=-2025/4,\quad
 L(-9)=-567/2,\quad L(-12)=-243/4.
\]
Using $p/(1+p)=p-p^2/(1+p)$ shows that the coefficients through $t^4$
cancel exactly for \eqref{prof:asymptotic-polynomial}.

The parameter interval \eqref{prof:parameter-displacement} is contained
in $[.40068272185,.40068272187]$.  Substituting this rational interval
in \eqref{prof:asymptotic-polynomial} and its $d$-derivative gives
\[
\begin{array}{c|rrrr}
 j&1&2&3&4\\ \hline
 |d_j|<&122&54200&12500000&5100000000\\
 |\partial_d d_j|<&304&100000&65000000&100000000000.
\end{array}
\]
These interval bounds use just the endpoint operations described in
Lemma~\ref{prof:root-bounds}; the coefficient formulas are polynomials.
The first row in particular implies
\begin{equation}\label{prof:tail-term-bounds}
 |d_1|56^{-3}<.00070,\quad |d_2|56^{-6}<1.8\cdot10^{-6},\quad
 |d_3|56^{-9}<8.1\cdot10^{-9},\quad |d_4|56^{-12}<3.2\cdot10^{-11}.
\end{equation}
Here is a finite positive-coefficient way to check the residual bound.
Let $b=(.00070,1.8\cdot10^{-6},8.1\cdot10^{-9},3.2\cdot10^{-11})$
and $B_*(z)=\sum_{i=1}^4b_iz^i$.  If $T_*=\sum b_i$, the sums of the
coefficients of degrees greater than four in its second, third, and fourth
powers are, respectively,
\[
 T_*^2-b_1^2-2b_1b_2-2b_1b_3-b_2^2,\qquad
 T_*^3-b_1^3-3b_1^2b_2,\qquad T_*^4-b_1^4.
\]
Substitution of the displayed four rational numbers gives
\begin{equation}\label{prof:residual-arithmetic}
 \sum_{j>4}[z^j]B_*(z)^2<8\cdot10^{-14},\qquad
 \sum_{j>4}[z^j]B_*(z)^3<2\cdot10^{-14},\qquad
 \sum_{j>4}[z^j]B_*(z)^4<3\cdot10^{-15}.
\end{equation}
Since $B_*(1)<.001$, all powers at least five contribute at most
$.001^5/(1-.001)$.  The remaining linear term contributes at most
$(243/4)(3.2\cdot10^{-11})/56^3$.  Thus the residual is bounded by
\[
 8\cdot10^{-14}+2\cdot10^{-14}+3\cdot10^{-15}
 +\frac{.001^5}{1-.001}
 +\frac{(243/4)(3.2\cdot10^{-11})}{56^3}
 <3\cdot10^{-13}.
\]
These bounds hold for all $s\ge56$ because every remaining absolute
monomial decreases there.  The residual is a rational function of
$s^{-3}$, analytic at zero with vanishing constant term, which also
proves the assertion about its derivatives.  For the parameter estimate,
the factors $3j/56$ and $3j(3j+1)/56^2$ are at most one for $1\le j\le4$.
Thus the second row of the coefficient table gives
\[
 \norm{\partial_dP_{\rm a}}_\infty
 \le\frac{304}{56^3}+\frac{100000}{56^6}
       +\frac{65000000}{56^9}+\frac{100000000000}{56^{12}}
 <.001735<.01.
\]
The mean value theorem yields \eqref{prof:tail-parameter-bound}.
\end{proof}

\begin{lemma}[Tail solutions]\label{prof:tail}
For each $d$ arising in \eqref{prof:parameters} and each real
$|q|\le .26R$, there is a solution $p_{\rm tail}$ of \eqref{prof:ode}
on $[56,\infty)$ such that
\[
 |p_{\rm tail}|<.001,
 \qquad q_{\rm tail}=q,\qquad |\xi_{\rm tail}|<.25R.
\]
Its first three jet entries tend to zero at infinity.  The solution,
as a bounded jet, and its initial coordinates depend continuously on $q,d$.
\end{lemma}

\begin{proof}
Let $b_i=D_0'(\lambda_i)^{-1}$ and let $j$ be a bounded continuous
real function on $[S_0,\infty)$.  Define a vector $\mathcal V=(V,V',V'')^t$
by the convergent integrals
\begin{equation}\label{prof:green-formula}
 \begin{split}
 \mathcal V(s)={}&q\e^{\lambda_0(s-S_0)}v_0
  +b_0v_0\int_{S_0}^s\e^{\lambda_0(s-u)}j(u)\,\dd u\\
 &-\sum_{i=1}^2 b_iv_i\int_s^\infty\e^{\lambda_i(s-u)}j(u)\,\dd u.
 \end{split}
\end{equation}
Indeed this vector solves the companion system for
$D_0(\partial_s)V=j$: the identities
$\sum b_i=\sum\lambda_ib_i=0$ and $\sum\lambda_i^2b_i=1$ follow
by partial fractions of $1/D_0$.  They show both the derivative relations
between its entries and the source in the last component.  The vector is
real by complex conjugation.  It is the unique bounded solution with the
prescribed real-mode coordinate at $S_0$, since a bounded homogeneous
solution has no growing nonreal modes.

For $0\le i\le2$, Lemma~\ref{prof:root-bounds} gives
\begin{equation}\label{prof:green-bound}
 \norm{V^{(i)}}_\infty\le |\lambda_0|^i|q|+M_i\norm j_\infty.
\end{equation}
In particular, on the closed ball $\norm j_\infty\le .3R$, the first
entry satisfies
$\norm V_\infty\le(.26+1.96\cdot.3)R=.848R$.
Together with \eqref{prof:tail-term-bounds}, this ensures
$|p_{\rm a}+V|<.001$ throughout the half-line.

Write $N(p)=p^2/(1+p)$ and $e_{\rm a}=\mathcal E_d(p_{\rm a})$.
Then $p=p_{\rm a}+V$ solves \eqref{prof:ode} exactly when
\begin{equation}\label{prof:tail-map}
 \begin{split}
 j=\mathcal T_{q,d}j:={}&-e_{\rm a}
 -\left(\frac{30}{s}-\frac{15}{S_0}\right)V''
 -\left(\frac B{s^2}-\frac B{2S_0^2}\right)V'\\
 &-\left(\frac C{s^3}-\frac C{2S_0^3}\right)V
 +N(p_{\rm a}+V)-N(p_{\rm a}).
 \end{split}
\end{equation}
For $|p|\le.001$,
\[
 |N'(p)|=|1-(1+p)^{-2}|<.0021.
\]
The three coefficient differences in \eqref{prof:tail-map} have absolute
bounds $15/S_0$, $B/(2S_0^2)$ and $|C|/(2S_0^3)$ on the entire half-line.
It follows from \eqref{prof:green-bound} that the Lipschitz constant of
$\mathcal T_{q,d}$ on the ball is at most
\begin{equation}\label{prof:tail-contraction-constant}
 \frac{15}{56}(1.88)+\frac{B}{2\cdot56^2}(1.92)
 +\left(\frac{|C|}{2\cdot56^3}+.0021\right)(1.96)<.57.
\end{equation}
The corresponding coefficient of the homogeneous term is bounded by
\begin{equation}\label{prof:tail-homogeneous-constant}
 \frac{15}{56}(1.08659)^2+\frac{B}{2\cdot56^2}(1.08659)
 +\frac{|C|}{2\cdot56^3}+.0021<.36.
\end{equation}
Consequently
\[
 \norm{\mathcal T_{q,d}j}_\infty
 \le3\cdot10^{-13}+.36|q|+.57\norm j_\infty
 \le(.0006+.36\cdot.26+.57\cdot.3)R=.2652R<.3R.
\]
The contraction theorem gives a unique fixed point in this ball.
Equation~\eqref{prof:green-formula} then gives a $C^3$ solution, and
\eqref{prof:ode} bootstraps it to a smooth solution for $s>0$.
Its nonreal initial coordinate satisfies
\[
 \xi_{\rm tail}=-b_1\int_{S_0}^\infty
                  \e^{\lambda_1(S_0-u)}j(u)\,\dd u,
 \qquad |\xi_{\rm tail}|<.82(.3R)=.246R<.25R.
\]

All coefficients in \eqref{prof:tail-map} depend continuously on $d$
in the relevant uniform norms, and the homogeneous vector depends
continuously on $q$.  The common contraction constant therefore gives
continuous dependence of the fixed point, the bounded jet, and its
initial coordinates.  Finally, exponential convolutions of continuous
functions tending to zero again tend to zero.  Thus
\eqref{prof:tail-map} preserves the closed subspace of such functions;
the residual tends to zero and all its other coefficients are bounded.
Iteration from $j=0$ shows that the fixed point belongs to this subspace.
Formula~\eqref{prof:green-formula} proves $p,p',p''\to0$.
\end{proof}

\subsection{Matching the two solutions}

For a parameter $z_*\in\mathcal Z$, apply Lemma~\ref{prof:tail} using
the value $d(z_*)$ and the coordinate $q_{\rm shot}(z_*)$ supplied by
the finite shot.  This is allowed by $|q_{\rm shot}|\le .25R<.26R$.
Define a map on the closed square by
\[
 \mathcal B(z_*)=z_*
 -\bigl(\Re(\xi_{\rm shot}-\xi_{\rm tail}),
        \Im(\xi_{\rm shot}-\xi_{\rm tail})\bigr).
\]
It is continuous by the continuous dependence established above and
satisfies
\[
 \norm{\mathcal B(z_*)}_\infty\le(.13+.25)R=.38R<R.
\]
Brouwer's fixed point theorem gives $z_*\in\mathcal Z$ with
$\xi_{\rm shot}=\xi_{\rm tail}$.  The real coordinate and the parameter
$d$ already agree.  The three jet entries in \eqref{prof:coordinates}
therefore agree, and ordinary uniqueness for \eqref{prof:ode} joins
the shot and tail to one solution on $[0,\infty)$.

The finite shot has $1+p>0$, and the tail has $|p|<.001$.
Compactness of the finite interval and the uniform tail bound imply
that $1+p$ is bounded and bounded away from zero globally.  Since
$1+d>0$, the function $G$ in \eqref{prof:change-variables} has the same
properties.  Near zero,
\[
 s^{3/2}=\frac{\sqrt{27}}{4c}y,
 \qquad G(y)=1+\frac{\sqrt{27}}4ky+O(y^2).
\]
It is therefore analytic there, with
$G'(0)=k\sqrt{27}/4<0$ by \eqref{prof:parameter-displacement}.
This proves all assertions of Theorem~\ref{prof:profile} except for the
differentiable expansion, which we now establish independently of the
finite shooting estimates.

\subsection{Differentiable asymptotics at every order}

The distinction between a bounded remainder and a remainder in a prescribed
polynomial weight matters here.  The following argument first proves
uniqueness among bounded solutions and then identifies the actual
remainder with the weighted solution.

\begin{lemma}[Weighted remainder]\label{prof:weighted-remainder}
Let $p$ be the matched solution above.  There are uniquely determined
coefficients $d_j$, extending \eqref{prof:asymptotic-polynomial}, such that
for every $J,\ell\ge0$,
\begin{equation}\label{prof:s-expansion}
 \frac{\dd^\ell}{\dd s^\ell}
 \left(p(s)-\sum_{j=1}^{J}d_js^{-3j}\right)
 =O\bigl(s^{-3(J+1)-\ell}\bigr).
\end{equation}
\end{lemma}

\begin{proof}
The limits $p,p',p''\to0$ and \eqref{prof:ode} imply that every derivative
of $p$ is bounded on a sufficiently late half-line.  To see this without
assuming any asymptotic expansion, solve \eqref{prof:ode} for $p'''$.
Its right-hand side is bounded, since $p$ is separated from $-1$.
Successively differentiating this formula proves the claim by induction:
the coefficient functions $s^{-j}$ have bounded derivatives and all
derivatives of $p\mapsto p/(1+p)$ are bounded on the range in question.

Formal substitution of $p=\sum_{j\ge1}d_jt^j$, $t=s^{-3}$, in
\eqref{prof:ode} determines the coefficients successively.  Indeed the
coefficient of $t^j$ contains the new coefficient $d_j$ with coefficient
one from $p/(1+p)$; every other term involves only previous coefficients.
The source $-Cd/s^3$ determines $d_1=Cd$.  For
$p_J=\sum_{j=1}^Jd_js^{-3j}$, including $p_0=0$, one consequently has
\begin{equation}\label{prof:formal-residual}
 e_J:=\mathcal E_d(p_J),\qquad
 e_J^{(\ell)}(s)=O\bigl(s^{-N-\ell}\bigr),
 \quad N=3(J+1),\quad \ell\ge0.
\end{equation}
This follows directly from analyticity of the finite expression in
$s^{-3}$ near zero; no convergence of the infinite formal series is asserted.

Let $W=p-p_J$.  The exact divided-difference identity
\[
 \frac p{1+p}-\frac{p_J}{1+p_J}
 =\frac{W}{(1+p)(1+p_J)}
\]
gives the linear equation
\begin{equation}\label{prof:remainder-equation}
 W'''+W=-e_J-a_2W''-a_1W'-a_0W,
\end{equation}
where
\[
 a_2=\frac{30}{s},\qquad a_1=\frac B{s^2},\qquad
 a_0=\frac C{s^3}+\frac1{(1+p)(1+p_J)}-1.
\]
These are now treated as given coefficient functions.  Their absolute
sum tends to zero because $p,p_J\to0$, and all their derivatives are
bounded on a sufficiently late half-line.

The polynomial $\lambda^3+1$ has roots $-1$ and
$(1\pm i\sqrt3)/2$.  The analogue of \eqref{prof:green-formula} on
$[S,\infty)$, with these roots, writes the bounded solution of
$W'''+W=j$ as a stable homogeneous vector plus a convolution vector
$\mathcal K_Sj$.  Every root has modulus one and
$|(\lambda^3+1)'|=3$ at the roots.  Hence
\begin{equation}\label{prof:weighted-green}
 \max_{0\le i\le2}\norm{(\mathcal K_Sj)_i}_\infty
 \le\frac53\norm j_\infty.
\end{equation}
It also acts in the Banach space with norm
$\norm j_N=\sup_{s\ge S}s^N|j(s)|$.  For $S\le u\le s$,
\[
 (s/u)^N\le\exp\bigl(N(s-u)/S\bigr),
\]
while for $u\ge s$ this ratio is at most one.  If $S\ge2N$, the
forward kernel is therefore bounded by an exponential with decay rate
$1/2$, and the two backward kernels already have decay rate $1/2$.
It follows that
\begin{equation}\label{prof:weighted-green-2}
 \max_{0\le i\le2}\sup_{s\ge S}s^N| (\mathcal K_Sj)_i(s)|
 \le2\norm j_N.
\end{equation}
The stable homogeneous vector $H(s)=q_S\e^{-(s-S)}(1,-1,1)^t$
belongs to both spaces.  Here $q_S$ is chosen to be the stable coordinate
of the actual jet $(W,W',W'')(S)$.

Choose $S\ge2N$ so late that
$2\sup_{s\ge S}(|a_0|+|a_1|+|a_2|)<1$.
The affine equation for the source,
\begin{equation}\label{prof:weighted-source}
 j=-e_J-\sum_{i=0}^2a_i\bigl(H_i+(\mathcal K_Sj)_i\bigr),
\end{equation}
is a contraction both in the bounded norm and in the norm $\norm\cdot_N$,
by \eqref{prof:weighted-green}--\eqref{prof:weighted-green-2}.
It has a unique solution in each space.  The weighted solution is bounded,
so uniqueness in the bounded space identifies the two solutions.
On the other hand, the actual remainder has a bounded jet and bounded
source $j=W'''+W$.  The growing-mode integrals from infinity are forced
by boundedness, so its jet has the representation
$H+\mathcal K_Sj$ with the chosen $q_S$.  It satisfies
\eqref{prof:weighted-source} and hence equals the weighted solution.
We have proved
\begin{equation}\label{prof:preliminary-remainder}
 W,W',W''=O(s^{-N}).
\end{equation}

Equation~\eqref{prof:remainder-equation} gives $W'''=O(s^{-N})$.
Differentiating it repeatedly, using the bounded coefficient derivatives
and \eqref{prof:formal-residual}, proves
$W^{(\ell)}=O(s^{-N})$ for every fixed $\ell$.
To obtain the extra decay for each derivative in
\eqref{prof:s-expansion}, fix $J,\ell$ and apply this last conclusion to
the deeper truncation $p_K$, where $K\ge J+\ell$.  Then
\[
 (p-p_J)^{(\ell)}
 =(p-p_K)^{(\ell)}+
   \sum_{j=J+1}^{K}d_j(s^{-3j})^{(\ell)}
 =O\bigl(s^{-3(J+1)-\ell}\bigr).
\]
The starting point $S$ and constants may depend on the finite orders
$J,\ell$, which is exactly what the asserted asymptotic statement requires.
\end{proof}

\begin{proof}[Completion of the proof of Theorem~\ref{prof:profile}]
Finally $s^{-3}=(16/27)c^2y^{-2}$, so
\eqref{prof:s-expansion} and the chain rule give
\eqref{prof:expansion}, with
$g_j=c\,d_j((16/27)c^2)^j$.  Repeated differentiation of
$u\mapsto G(u^{-1/2})$ expresses each derivative as a finite sum of
powers of $y$ times derivatives of $G$.  Subtracting sufficiently many
terms in \eqref{prof:expansion} shows that these derivatives have their
Taylor limits at $u=0$.  The resulting extension is $C^\infty$.
The origin expansion gives the smooth extension of $h$, and the
calculation in Section~\ref{geo:core-section} proves the flat
equation in Theorem~\ref{prof:profile}, also at $x=0$ by smoothness.
\end{proof}
